% !PRINTED 341
\weberpage{341}
% running head: §. 91. | Reality. | 341
%<<CONT

\weq{18}{\begin{array}{ccc}
(0, 1, -1), & (-\varepsilon, 0, 1), & (1, -\varepsilon^2, 0) \\
(-1, 0, 1), & (1, -\varepsilon, 0), & (0, 1, -\varepsilon^2) \\
(1, -1, 0), & (0, 1, -\varepsilon), & (-\varepsilon^2, 0, 1),
\end{array}}
\weq{19}{\begin{array}{ccc}
(0, 1, -1), & (-\varepsilon^2, 0, 1), & (1, -\varepsilon, 0) \\
(-1, 0, 1), & (1, -\varepsilon^2, 0), & (0, 1, -\varepsilon) \\
(1, -1, 0), & (0, 1, -\varepsilon^2), & (-\varepsilon, 0, 1).
\end{array}}

Of the straight lines (16), each contains one real and two
conjugate imaginary points, and all three lines are real. Of the
lines (17), the first contains the three real points and is real,
while the other two are conjugate imaginary. In (18) and (19)
all three lines are imaginary, and indeed the lines of (18) are
conjugate to the lines of (19).

A still clearer notation is the following. If
$\xi$ and $\eta$ each run through a complete system of residues with respect to the modulus 3,
say $0, 1, 2$, and two numbers congruent with respect to the modulus 3
are regarded as not distinct, then there are nine
combinations $(\xi, \eta)$, which we can use as a designation for the nine
points of inflection. If we stipulate that $(\xi, \eta)$ and
$(\eta, \xi)$ denote conjugate imaginary points of inflection, and hence $(0, 0)$, $(1, 1)$, $(2, 2)$
the three real ones, then from (16) to (19) we obtain the following
 table, in which again the points of inflection standing in one row
 lie on a straight line:
\weq{20}{\begin{array}{ccc@{\qquad}ccc}
(0,0) & (1,2) & (2,1) & (0,0) & (1,1) & (2,2) \\
(1,1) & (2,0) & (0,2) & (1,2) & (2,0) & (0,1) \\
(2,2) & (0,1) & (1,0) & (2,1) & (0,2) & (1,0) \\[4pt]
(0,0) & (2,0) & (1,0) & (0,0) & (0,2) & (0,1) \\
(1,1) & (0,1) & (2,1) & (1,1) & (1,0) & (1,2) \\
(2,2) & (1,2) & (0,2) & (2,2) & (2,1) & (2,0).
\end{array}}

One sees from this that three points of inflection $(\xi_1, \eta_1)$, $(\xi_2, \eta_2)$,
$(\xi_3, \eta_3)$ lie on a straight line then and only then, when
the two congruences
\weq{21}{\begin{aligned}
\xi_1+\xi_2+\xi_3 &\equiv 0 \\
\eta_1+\eta_2+\eta_3 &\equiv 0
\end{aligned} \quad (\mathrm{mod}\ 3)}
are satisfied.
